\subsection{关于紧可度量化拓扑空间的补充结果}

我们给二元素集 \(\{0,1\}\) 赋以离散拓扑，给集合 \(\{0,1\}^{N}\) 赋以乘积拓扑。拓扑空间 \(\{0,1\}^{N}\) 是紧的（TG, I, p. 63, 定理 3）、可度量化的（TG, IX, p. 15, 推论 2）、非空的、全不连通的（TG, I, p. 84, 命题 10），且没有孤立点。

命题 14. --- 每个紧、可度量化、非空、全不连通且无孤立点的拓扑空间都同胚于 \(\{0,1\}^{N}\)。

设 X 为这样的拓扑空间。给它赋以与它的拓扑相容的距离。由于空间 X 紧，它对这个距离是完备的（TG, II, p. 27, 定理 1）。

引理 3. --- 设 \(\varepsilon\) 为实数 \(>0\)。存在整数 \(m \geqslant 1\)，使得对每个整数 \(n \geqslant m\)，X 都有一个由 \(n\) 个直径 \(\leqslant \varepsilon\) 的非空开闭集组成的分割。

X 的每一点都有直径 \(\leqslant\varepsilon\) 的开闭邻域（TG, II, p. 32, 命题 6 的推论）。由于空间 X 紧，它有一个由这样的集合组成的有限覆盖。选取其中之一 \((\mathrm{U}_{i})_{1\leqslant i\leqslant m}\)，使 m 最小。由于 X 非空，有 \(m\geqslant1\)。对 \(1\leqslant i\leqslant m\)，用 \(V_{i}\) 表示 \(U_{i}\) 与

\(X\setminus U_k\)（\(k < i\)）的交。于是 \((V_i)_{1 \leqslant i \leqslant m}\) 是 X 的一个分割，由 m 个直径 \(\leqslant \varepsilon\) 的非空开闭集组成。

由于 X 没有孤立点，X 的每个非空开闭子集 V 都至少含两点。此外，由于 X 紧且全不连通，V 是两个不交的非空开闭子集的并（同前所引）。由此归纳地得出，对每个整数 \(n \geqslant m\)，X 都有一个由 n 个直径 \(\leqslant \varepsilon\) 的非空开闭集组成的分割。

现在来完成命题 14 的证明。X 的每个非空开闭子空间都是紧、全不连通、无孤立点的度量空间。于是引理 3 使我们可以归纳地构造有限集的序列 \((\mathrm{J}_{n})_{n\geqslant0}\)，以及对每个 \(n\geqslant0\)，一个映射 \(\varphi_{n}\)（它从集合 \(C_{n}=J_{0}\times\cdots\times J_{n}\) 映到 X 的直径 \(\leqslant2^{-n}\) 的非空开闭子集组成的集合中），使得：

\begin{enumerate}
\def\labelenumi{(\roman{enumi})}
\item
  对每个 \(n \geqslant 0\)，存在整数 \(m_n \geqslant 1\)，使得 \(\mathrm{Card}(\mathrm{J}_n) = 2^{m_n}\)；
\item
  族 \((\varphi_0(c))_{c\in \mathrm{C}_0}\) 是 \(X\) 的一个分割；
\item
  对每个 \(n \geqslant 0\) 及每个 \(c \in \mathrm{C}_n\)，族 \((\varphi_{n+1}(c,j))_{j \in \mathrm{J}_{n+1}}\) 是 \(\varphi_n(c)\) 的一个分割。
\end{enumerate}

用 \(p_n\) 表示从 \(C_{n+1}\) 到 \(C_n\) 上的典范投影。序列 \(C = (C_n, p_n)_{n \geqslant 0}\) 是一个筛（TG, IX, p. 63, 定义 8）。与这个筛相联系的拓扑空间（TG, IX, p. 63）等同于拓扑空间 J，即诸离散拓扑空间 \(J_n\) 的乘积。筛 C 与映射序列 \((\varphi_n)_{n \geqslant 0}\) 定义了度量空间 X 的一个严格筛分（TG, IX, p. 63 与 p. 64）。由这个筛分得到的映射 \(f: J \to X\) 是连续的双射（TG, IX, p. 65）。由于拓扑空间 J 是紧的（TG, I, p. 63, 定理 3）而 X 是分离的，\(f\) 是同胚（TG, I, p. 63, 推论 2）。由于 \(J_n\) 同胚于 \(\{0, 1\}^{m_n}\)（\(n \geqslant 0\)），J 同胚于 \(\{0, 1\}^N\)（TG, I, p. 25, 命题 2）。

例. --- 设 K 为康托尔三分集（TG, IV, p. 9, 例）。对每个 \(n \geqslant 0\)，令 \(J_{n} = \{0,1\}\)，并如下定义映射 \(K_{n}\)（从集合 \(C_{n} = J_{0} \times \cdots \times J_{n}\) 到 {[}0,1{]} 的闭区间组成的集合中）：令 \(K_{0}(0) = [0, \frac{1}{3}]\) 与 \(K_{0}(1) = [\frac{2}{3}, 1]\)；对每个 \(n \geqslant 0\) 与每个 \(c \in C_{n}\)，\(K_{n+1}(c,0)\) 与

\(K_{n+1}(c,1)\) 分别是 \(K_n(c)\) 的“左三分”与“右三分”。若 \(c = (j_0,j_1,\ldots,j_n) \in C_n\)，则 \(K_n(c)\) 是（同前所引）中记作 \(K_{n,p}\) 的区间，其中 \(p = 2^n j_0 + 2^{n-1}j_1 + \cdots + j_n + 1\)；它也是区间 \([a,a + \frac{1}{3^{n+1}}]\)，其中 \(a = 2(\frac{j_0}{3} + \frac{j_1}{3^2} + \cdots + \frac{j_n}{3^{n+1}})\)。对 \(n \geqslant 0\) 与 \(c \in C_n\)，令 \(\varphi_n(c) = K_n(c) \cap K\)。族 \((\varphi_n(c))_{c \in C_n}\) 是由闭集组成的 K 的一个分割。因此这些集合在 K 中也是开的；它们非空且直径为 \(\frac{1}{3^{n+1}}\)，因为区间 \(K_n(c)\) 的端点都属于 K。序列 \(C = (C_n, p_n)_{n \geqslant 0}\)（其中 \(p_n: C_{n+1} \to C_n\) 是典范投影 \((c,j) \mapsto c\)）是一个筛，与这个筛相联系的拓扑空间等同于 \(\{0,1\}^N\)。筛 C 与映射序列 \((\varphi_n)_{n \geqslant 0}\) 定义了度量空间 K 的一个严格筛分。由这个筛分得到的映射 \(f: \{0,1\}^N \to K\) 是同胚，由下列公式给出

\[
f ((j _ {n}) _ {n \geqslant 0}) = 2 \sum_ {n = 0} ^ {\infty} \frac {j _ {n}}{3 ^ {n + 1}}.
\]

推论. --- 设 X 为可度量化、紧且非空的拓扑空间。存在从 \(\{0,1\}^{N}\) 到 X 上的连续满射。

每个紧可度量化的拓扑空间都同胚于拓扑空间 \(I^{N}\) 的一个子空间，且该子空间必然是闭的（TG, IX, p. 18, 命题 12 与 p. 21, 命题 16）。设 A 为 \(I^{N}\) 的闭子空间，h 为从 A 到 X 上的同胚。

令 \(K = \{0,1\}^{N}\)，且对 \(\alpha = (a_n) \in K\)，令 \(f(\alpha) = \sum_{n=0}^{\infty} a_n 2^{-n-1}\)；我们有 \(f(\alpha) \in I\)。这样定义的映射 \(f: K \to I\) 是满射（TG, IV, p. 42）且连续。事实上，若 K 的两个元素 \(\alpha\) 与 \(\beta\) 的指标 \(< n\) 的坐标都相同，我们有 \(|f(\beta) - f(\alpha)| \leqslant 2^{-n}\)。用 \(g\) 表示映射 \((\alpha_n) \mapsto (f(\alpha_n))\)（从 \(K^N\) 到 \(I^N\) 中）；它是连续的满射。拓扑空间 \(K^N\) 是紧的（TG, I, p. 63, 定理 3）、可度量化的（TG, IX, p. 15, 推论 2）且全不连通的（TG, I, p. 84, 命题 10），它的闭子空间 \(\bar{g}^{1}(A)\) 亦然；后者是非空的，因为映射 \(g\) 与 \(h\) 都是满射。

于是，空间 \(\bar{g}^{1}(A) \times K\) 同胚于 \(\{0,1\}^{N}\)（命题 14），因为它是紧、可度量化、全不连通且无孤立点的拓扑空间，而映射 \((x,y) \mapsto h(g(x))\) 从 \(\bar{g}^{1}(A) \times K\) 到 X 中是连续且满射的。

